\chapter*{Acronyms}
\label{chap:acro}
\addcontentsline{toc}{chapter}{Acronyms}

\begingroup
\setlength{\itemsep}{0.25em}
\begin{acronym}[SHA-256]
  \acro{API}{Application Programming Interface}
  \acro{C2}{Command and Control (command-and-control server)}
  \acro{CLI}{Command Line Interface}
  \acro{DLL}{Dynamic Link Library}
  \acro{ETW}{Event Tracing for Windows}
  \acro{FR}{Functional Requirement}
  \acro{GUI}{Graphical User Interface}
  \acro{HTTP}{Hypertext Transfer Protocol}
  \acro{IL}{Intermediate Language (.NET)}
  \acro{JSON}{JavaScript Object Notation}
  \acro{LLM}{Large Language Model}
  \acro{NFR}{Non-Functional Requirement}
  \acro{NDJSON}{Newline-Delimited JSON (event streaming)}
  \acro{PE}{Portable Executable}
  \acro{PsDirect}{PowerShell Direct (remote access to Hyper-V guest)}
  \acro{RAT}{Remote Access Trojan}
  \acro{REST}{Representational State Transfer}
  \acro{RQ}{Redis Queue}
  \acro{SQL}{Structured Query Language}
  \acro{UBI}{Universidade da Beira Interior}
  \acro{UUID}{Universally Unique Identifier}
  \acro{VM}{Virtual Machine}
  \acro{SHA-256}{Secure Hash Algorithm 256 bits (file integrity)}
  \acro{WPF}{Windows Presentation Foundation}
  \acro{YARA}{Yet Another Recursive Acronym (detection rule language)}

\end{acronym}
\endgroup
